\chapter{Self-Gravitating Gas}\label{ch:cosmology}

\small
\begin{quote} According to the \emph{Perfect Cosmological Principle}, 
the Universe on a large scale looks the same everywhere and at all times.
(Fred Hoyle)
\end{quote}
\begin{quote} The thermodynamic properties of self-gravitating
systems are still unclear...The tendency of increasing granularity with time
(in a self-gravitating system), so important to the structure and development 
of the Universe, seems to represent a fundamental principle. 
(Paul Davies \cite{davies})
\end{quote}
\begin{quote} 
GRAVITATION: The universal property of all material objects
to attract each other.\\
GRAVITON: A hypothetical particle which, when passing
to and fro in virtual form bewteen two masses, in large numbers,
is thought to mediate the gravitational attraction between them.
(Guillemin \cite{guillemin})
\end{quote}
\normalsize

\normalsize

\section{Euler's Equations with Gravitation}

The Euler equations including gravitational forces from the mass distribution
can be viewed as a basic (non-relativistic) \emph{cosmological model}.  
We thus consider an inviscid perfect gas 
enclosed in a volume $\Omega$ in $\bbR^3$ with boundary $\Gamma$
(or the whole of $\bbR^3$), over a time interval $I=(0,1]$, assuming 
that the gas is subject to a gravitational force $\nabla\varphi$, where
$\varphi$ is a \emph{gravitational potential} 
coupled (at each time instant) to the mass distribution $\rho$  
through Poisson's equation $\Delta\varphi =\rho$
with (for simplicity) $\varphi =0$ on $\Gamma$. This is a 
\emph{self-gravitating system} since the gravitational force depends
on the mass-distribution  

Replacing $\rho$ by $\Delta\varphi$, we are led to the following version
of Euler's equations with gravitation:  
Find $\hat u=(\varphi ,m,e)$ 
depending on $(x,t)\in Q\equiv \Omega\times I$ such that
\begin{equation}\label{eulergrav}
\begin{array}{rcl}
\Delta\dot\varphi +\nabla\cdot m &=&  0 \quad \mbox{in } Q, \\
\dot m +\nabla\cdot (mu) +\nabla p-\rho\nabla\varphi&=&  0 \quad \mbox{in } Q, \\
\dot e +\nabla\cdot (eu)+p\nabla\cdot u &=& 0  \quad \mbox{in } Q,\\ 
u\cdot n&=&0\quad \mbox{on } \Gamma\times I\\
\varphi &=&0\quad \mbox{on } \Gamma\times I\\
\hat u(\cdot ,0)&=&\hat u^0\quad \mbox{in } \Omega , 
\end{array} 
\end{equation}
where $u=\frac{m}{\Delta\varphi}$, 
$p=(\gamma -1)e$ with $\gamma >1$, and we normalize the gravitational
constant to unity.

\section{The Hen and the Egg}

An interesting aspect of this model with the gravitational
potential $\varphi$ as a primary variable and the mass distribution
$\rho$ as a derived quantity, is that it does not include
\emph{action at distance} in the same way as the conventional model
with $\varphi$ derived from $\rho$ through Poisson's equation 
$\Delta\varphi =\rho$. Instead we have to deal
with the new equation $\Delta\dot\varphi +\nabla\cdot m =0$,
still containing the Laplacian.

Considering the equation $\Delta\varphi =\rho$ as defining
$\varphi$ in terms of $\rho$ or vice versa, of course connects
to which comes first: the egg or the hen. 

Choosing (in an exam for example) 
to explain either how a hen can lay an egg
or how an egg can develop into a hen, you may go for the first
option as appearing to be (somewhat) simpler to deal with.
The idea would be that it would be advantageous to
choose the more complex object of the 
potential/hen rather than the simpler mass/egg, as 
the primary unknown which will come out by solving the equations,
(by some form of computation).
Thus the more complex object, the potential/hen will 
be ``created by computation'' in front of our eyes, thus relieving us
from \emph{explaining} how a mass/egg can create a potential/hen. 
But doing so we have to deal with 
the equation $\Delta\dot\varphi +\nabla\cdot m=0$
expressing ``mass conservation'' in a non-standard form,
which  admittedly seems to involve action at distance if we rewrite it in the form.
$\dot\varphi +(\Delta )^{-1}\nabla\cdot m=0$ for the purpose
of explicitly updating $\varphi$ by time-stepping. What can we do 
to allow time-stepping without inverting the Laplacian?

\section{Perturbed Mass Conservation} 

Well, we may consider
to perturb the equation for mass conservation
$\Delta\dot\varphi +\nabla\cdot m=0$ into: 
\[
\kappa\ddot\varphi -  \Delta\dot\varphi -\nabla\cdot m=0
\]
with $\kappa$ a small positive constant, now allowing time stepping
without inverting the Laplacian. Of course, you could regularize the
equation $\Delta\varphi =\rho$ similarly into  
$\kappa\dot\varphi -\Delta\varphi =\rho$ keeping mass conservation in the
standard form $\dot\rho +\nabla\cdot (\rho u)=0$. In any case,
(\ref{eulergrav}) with $\varphi$ instead of $\rho$ and a perturbed
equation for mass conservation, is a model without action at distance.


\section{Gravitons?}

Solving Poisson's equation $\Delta\varphi =\rho$
equation for the gravitational
potential $\varphi$ in terms of the density $\rho$, reflect that
(somehow) a mass distribution ``creates'' its own gravitational field
(seemingly infinitely fast). The physics of this creation of a hen out
of an egg in the form of
a gravitational potential/force from a mass distribution, is unknown: A
conjecture is that the gravitational field is 
established through an \emph{exchange} of certain hypothetical particles 
called \emph{``gravitons''}. But no gravitons have been detected despite
intense search.



The reverse question is how a hen can lay an egg 
or how a gravitational potential can ``create'' a mass.
A potential with a sharp peak of the form 
$\frac{1}{4\pi\vert x-\bar x\vert}$, would then represent a
unit mass at position $\bar x$.  
In this perspective a tendency of ``mass lumping'' from 
very small initial density
variations, is to be expected from the fact that application of the Laplacian
may enhances variations and thus create positive and negative peaks/masses
out of nothing, with masses of different equal signs attracting each other
and masses with different signs repelleing each other. Our Universe
with positive masses would thus have a twin Universe with negative masses.

In this model, it is thus the gravitational 
potential $\phi$, which is given initially
and which evolves in time (together with momentum and energy)
and from which the mass density $\rho$ is generated by $\Delta\phi =\rho$
in a local ``creation process''.
We here avoid action at distance and we also get a hint
to the possible nature of all the ``dark mass'' seemingly required to account 
for the observed gravitational forces on cosmic scales, as 
mollified potential peaks not strong enough to 
``create visible mass'', but yet having gradients and exerting
gravitational forces.

%The alternative model (\ref{cosmology2}) with $\delta >0$ is a wave equation 
%with finite speed of propagation depending on $\delta$. The proper choice
%of $\delta$ from physical point of view is unclear, while 
%we expect certain outputs to be insensitive to the choice of $\delta$.

\section{The 2nd Law with Gravitation}

The 2nd Law with gravitation takes the form
\begin{equation}
\dot K=W-D+ P,\quad \dot E=-W+D,\quad D>0,
\end{equation}
where 
\[
\begin{split}
P&=-\int_\Omega \rho\nabla\varphi\cdot u\, dx =+\int_\Omega \varphi\nabla\cdot m\, dx=-\int_\Omega\varphi\dot\rho\, dx= \dot\Phi,\\
\Phi &=-\frac{1}{2}\int_\Omega \varphi\Delta\varphi\,dx =\frac{1}{2}\int_\Omega\vert\nabla\varphi\vert^2\, dx.
\end{split}
\] 
By summation, we have
\[
\frac{d}{dt}(K+E -\Phi)=0
\]
stating that the total energy, the sum of
kinetic, heat and gravitational energy $K+E -\Phi$, is conserved.

\section{Irreversibility and Heat Death}

A natural scenario is to start with a hot compressed gas at rest, which is
let free to expand in a \emph{Big Bang} with $W-D-P>0$ setting the
gas in motion until $W-D-P<0$ and $K=0$, whereupon the gas
contracts by gravitation back to a hot compressed state in a \emph{Big Crunch}
allowing the process to start all over with a new Big Bang.

In each cycle of this process some kinetic energy would irreversibly
be transformed into heat energy, which ultimately could lead to 
a stationary state with $W=D$ and $P=0$, a \emph{heat death}. But the 
whole process may have many cycles....and so there may have been Worlds
and intelligent cultures including thermodynamics  
even before the Big Bang of the World we happen to live in...   

\section{A Self-Gravitating Gas in a Box}

To probe the thermodynamics of (\ref{eulergrav}) concretely we solve a
discrete form on a cubic grid filling a box $\Omega =[0,1]^3$, in conservation
form for the density $\rho$, the momentum $m=\rho u$ and the internal energy
density $e$, with the ideal-gas closure
\begin{equation}\label{closure}
p=(\gamma -1)e=\rho T,\qquad \gamma >1,
\end{equation}
so that $T=p/\rho$ is the temperature and $\gamma$ the adiabatic index. The
self-gravitation is supplied through the \emph{mean-subtracted} Poisson coupling
\begin{equation}\label{jeansswindle}
\Delta\varphi = G\,(\rho -\bar\rho ),\qquad
\bar\rho =\frac{1}{\vert\Omega\vert}\int_\Omega\rho\, dx,
\end{equation}
with $\varphi =0$ on $\Gamma$, where $G>0$ is the \emph{gravitational coupling}
(playing the role of $4\pi G_{\mathrm{Newton}}$) and $\bar\rho$ the mean density.

The subtraction of $\bar\rho$ is essential and is the classical
\emph{Jeans swindle}. A strictly uniform gas in a \emph{finite} box is not
force-free if $\rho$ itself sources $\varphi$: with $\varphi =0$ on $\Gamma$ the
solution of $\Delta\varphi =G\rho$ bulges in the interior, so a net inward force
develops everywhere and the whole box \emph{breathes} in and out --- an artifact
of the boundary rather than of any structure in the gas. Sourcing $\varphi$ from
the density \emph{contrast} $\rho -\bar\rho$ makes the uniform background exactly
force-free, so that only departures from uniformity gravitate. This is precisely
what one wants of a cosmological model: it is the \emph{granularity}, not the
mean, that drives the dynamics.

The model so written has exactly \emph{two} control parameters: the
gravitational coupling $G$ and the adiabatic index $\gamma$. Everything in the
competition between collapse and dispersal is decided by these two, as we now
show.

\section{Pressure versus Gravitation: the Jeans Criterion}

Linearize (\ref{eulergrav}), (\ref{closure}), (\ref{jeansswindle}) about a
uniform state $\rho =\bar\rho$, $u=0$, $T=\bar T$ at rest. Writing the density
perturbation as $\delta\rho$ and eliminating $m$ and $\varphi$ one obtains the
wave equation
\begin{equation}\label{waveeq}
\ddot{\delta\rho}=c_s^2\,\Delta(\delta\rho)+G\bar\rho\,\delta\rho,
\qquad c_s^2=\gamma\,\frac{\bar p}{\bar\rho}=\gamma\bar T,
\end{equation}
where $c_s$ is the adiabatic sound speed. For a plane wave
$\delta\rho\propto\exp(i(k\cdot x-\omega t))$ this gives the
\emph{Jeans dispersion relation}
\begin{equation}\label{dispersion}
\omega^2=c_s^2k^2-G\bar\rho=\gamma\bar T\,k^2-G\bar\rho .
\end{equation}
The first term, from pressure, is \emph{stabilizing} (it grows with $\gamma$);
the second, from gravitation, is \emph{destabilizing} (it grows with $G$). A
perturbation grows --- the gas \emph{collapses} --- precisely when
$\omega^2<0$, i.e.\ when
\begin{equation}\label{jeanscrit}
G\bar\rho>\gamma\bar T\,k^2
\qquad\Longleftrightarrow\qquad
\lambda>\lambda_J=2\pi\sqrt{\frac{\gamma\bar T}{G\bar\rho}}\, ,
\end{equation}
$\lambda_J$ being the \emph{Jeans length}. Below the Jeans length pressure wins
and the perturbation merely oscillates as a sound wave ($\omega^2>0$); above it
gravity wins and the perturbation grows.

The decisive observation is that the onset of collapse depends on $G$ and
$\gamma$ \emph{only through the ratio} $G/\gamma$: at a fixed scale $k$ and
fixed background $(\bar\rho ,\bar T)$, doubling the coupling $G$ has exactly the
same effect as halving the stiffness $\gamma$. The dimensionless number that
decides everything is
\begin{equation}\label{jeansnumber}
J=\frac{G\bar\rho}{\gamma\bar T\,k^2},
\qquad\text{collapse}\iff J>1 .
\end{equation}

In a box of side $L$ the longest available wavelength is $\lambda\approx L$, so
structure forms \emph{at all} only if $\lambda_J\lesssim L$, that is
\begin{equation}\label{boxthreshold}
G\gtrsim (2\pi)^2\,\frac{\gamma\bar T}{\bar\rho\,L^2}.
\end{equation}
For $\gamma =2$ and $\bar T=\bar\rho =L=1$ this requires $G\gtrsim 79$; for a
seed of size $\sim L/3$ the threshold is about nine times larger. Below it the
box is Jeans-stable at \emph{every} scale and any initial lump simply pulsates.
This is why a modest coupling produces only an oscillation and never a collapse:
$G$ has not been pushed past the Jeans threshold set by $\gamma$.

\section{Collapse to a Stable Configuration: the Role of $\gamma$}

The Jeans criterion (\ref{jeanscrit}) decides whether collapse \emph{begins}.
Whether the collapse \emph{ends} in a stable, finite lump --- rather than
running away to a singular spike --- is a separate question, and it is decided
by $\gamma$ \emph{alone}, independently of $G$.

Consider a self-gravitating polytropic cloud of fixed mass $M$ and radius $R$.
Under adiabatic compression $p\propto\rho^\gamma$, so the internal energy
density scales as $e\propto\rho^\gamma\propto R^{-3\gamma}$ and the total
internal (thermal) energy as
\[
U=\int_\Omega e\,dx \;\propto\; R^{-3\gamma}\cdot R^3 \;=\; R^{-3(\gamma -1)},
\]
while the gravitational self-energy scales as
\[
W \;\propto\; -\frac{GM^2}{R}\;\propto\; -R^{-1}.
\]
The total energy is therefore $E(R)=A\,R^{-3(\gamma -1)}-B\,R^{-1}$ with
$A,B>0$. As $R\to 0$ the thermal term dominates --- and hence $E$ possesses a
\emph{minimum} at a finite radius, a stable hydrostatic equilibrium --- if and
only if it stiffens faster than gravity, $3(\gamma -1)>1$, that is
\begin{equation}\label{fourthirds}
\boxed{\;\gamma>\tfrac43\;}
\end{equation}
the classical \emph{Chandrasekhar index}. The three regimes are:
\begin{itemize}
\item $\gamma>\tfrac43$: pressure wins under compression; collapse is
\emph{arrested} at a finite quiescent lump, possibly after damped oscillations.
\item $\gamma<\tfrac43$: gravity wins at every compression; collapse
\emph{runs away}, no equilibrium exists.
\item $\gamma=\tfrac43$: the marginal case (governing, e.g., white-dwarf
stability).
\end{itemize}
Thus the two parameters play complementary roles: $G$ (through (\ref{jeanscrit}))
selects \emph{which scales collapse and how fast}, while $\gamma$ (through
(\ref{fourthirds})) decides \emph{whether the collapse can be stopped}. To reach
a stable bound configuration one needs \emph{both} a coupling $G$ large enough
to cross the Jeans threshold \emph{and} a stiffness $\gamma>\tfrac43$ so that the
endpoint is an equilibrium rather than a singularity.

\section{Breathing, Sloshing and Virialization}

These two criteria account for the phenomenology seen in the simulations.

\emph{Breathing.} A uniform gas in which $\rho$ (rather than $\rho-\bar\rho$)
sources $\varphi$ is not in equilibrium in a finite box: it contracts toward the
centre, pressure builds, it rebounds, and the box breathes globally while its
density stays nearly uniform. The Jeans swindle (\ref{jeansswindle}) removes
this boundary artifact, leaving only the genuine gravitational response to
density contrast.

\emph{Sloshing.} A pair of overdensities below the Jeans threshold
($J<1$ at their separation) does not collapse but exchanges mass back and forth
quasi-periodically between the two centres. This is a \emph{virial oscillation}
in the double-well gravitational potential: gas falling toward one centre
overshoots on its pressure and inertia, sloshes to the other, and returns. With
$\gamma>\tfrac43$ and no dissipation the oscillation is undamped, kinetic and
gravitational energy being traded reversibly, $\dot K=P=\dot\Phi$, exactly as in
the energy balance of \S\,\textit{The 2nd Law with Gravitation}.

\emph{Virialization.} To settle into a quiescent lump the gas must \emph{shed}
its infall energy: a stable bound configuration requires dissipation $D>0$ to
convert the ordered kinetic energy of collapse into heat. The entropy so
produced is the irreversible price of structure formation --- the mechanism
behind Davies' ``increasing granularity'' (\S\,\textit{Cosmology}) and behind
the heat death of \S\,\textit{Irreversibility and Heat Death}. Gravitational
structure forms by the conversion, through $D$, of gravitational potential
energy first into bulk motion and then, irreversibly, into heat.

\begin{remark}
In the numerical model the closure is written as $p=\gamma_s\,e$, so that the
adiabatic index is $\gamma=1+\gamma_s$; the stability boundary (\ref{fourthirds})
then reads $\gamma_s>\tfrac13$, and the coupling $G$ is exactly the
$4\pi G_{\mathrm{Newton}}$ of (\ref{jeansswindle}). The two sliders $(G,\gamma_s)$
of the interactive model are precisely the two parameters $(G,\gamma)$ of the
analysis above.
\end{remark}

\section{RealNSE Simulations of a Self-Gravitating Gas}

We give below some results from RealNSE solutions of (\ref{eulergrav})
showing the development mass granularity from
uniform mass distributions under expansion.

\begin{figure}[bhpt] 
\centerline{ 
\includegraphics[height=8cm]{ambsthermopict/cosmology.eps} 
} 
\caption{Cosmology theories}
\label{cosmologytheory} 
\end{figure}



\begin{figure}[bhpt]
\centerline{
\includegraphics[height=8cm]{ambsthermopict/cosmology2.eps}
}
\caption{Big Bang expansion}
\label{bigbang2}
\end{figure}

\section{The Basic Case: Gentle Collapse to a Pressure-Supported Core}\label{sec:selfgrav_basic}

The basic phenomenon of a self-gravitating gas is exhibited by the
browser-runnable 3D simulator \texttt{selfgrav\_collapse\_3d\_gpu.html} (in the
\texttt{Claes542/RealMolecule} repository). An \emph{isolated} cloud
($\Delta\varphi = G\rho$ with $\varphi = 0$ on the box, $p = \gamma\rho T$) is
started as a \emph{warm, mildly super-Jeans} over-density at rest --- warm enough
that the infall stays subsonic, so the dynamics is gentle and fully resolvable
on the $100^3$ grid. With $\gamma = 1$ (so $\Gamma = 1+\gamma = 2 > 4/3$, the
\emph{stable} regime of the role-of-$\gamma$ analysis above) the cloud runs
through the canonical sequence:

\begin{enumerate}\itemsep1pt
\item \textbf{Collapse.} Gravity exceeds the pressure gradient and the cloud
contracts, the centre running up in density and --- by compression and the
dissipation $\nu|\nabla u|^2$ --- in temperature.
\item \textbf{Bounce.} The rising central pressure halts the infall, but the
inward momentum overshoots hydrostatic balance, so the core rebounds.
\item \textbf{Ring-down.} The overshoot oscillates and is \emph{damped} by the
dissipation: each bounce carries steep velocity gradients whose kinetic energy
is thermalised, so the oscillation decays --- violent relaxation in miniature.
\item \textbf{Pressure-supported core.} The cloud settles into a
quasi-hydrostatic core, a Newtonian ``star''.
\end{enumerate}

\paragraph{Reading the force balance.}
Among the centre-row line plots the simulator draws the two terms of the
momentum equation directly: \emph{white} $|\nabla p|$ (pressure) and \emph{red}
$\rho|\nabla\varphi|$ (gravity). During collapse the red curve runs above the
white one (gravity winning); once the core settles the two curves \emph{lie on
top of each other}, which is hydrostatic balance $\nabla p = -\rho\nabla\varphi$
recovered. In the settled state the white curve sits a hair \emph{below} the
red: the pressure gradient is marginally weaker than gravity, a small net inward
force corresponding to the slow quasi-static contraction of the core (part
physical, part the residual draining of pressure support by the numerical
viscosity). The gap is a live measure of how close the discrete solution is to
exact hydrostatic balance.

\paragraph{Breakdown and recovery.}
Where the flow drives gradients sharper than the mesh can hold, the smooth
solution \emph{breaks down}; the dissipation thermalises the broken-down
small-scale energy and a smoother mean solution \emph{re-emerges}. This
breakdown-and-recovery is the 2nd Law of Chapter~\ref{ch:essence} operating in
the self-gravitating setting: the pointwise fine structure at the breakdown is
lost irreversibly, but the outputs --- the density profile, the temperature, the
force balance --- recover and carry on. The analytical counterpart is the
$\gamma = 1$ polytrope of the Lane--Emden family, $\rho \propto \sin(kr)/(kr)$,
a finite stable sphere; the simulator finds the same equilibrium dynamically,
as the attractor of the damped collapse.

\paragraph{Why warm, and the limit of the model.}
A \emph{cold} (near-pressureless) start makes the infall supersonic and the
bounce near-singular --- a core a few cells wide with an enormous sound speed,
which no fixed grid resolves; that is itself the correct statement that the
continuum is heading for a singularity where new physics must enter. For
$\gamma < 1/3$ ($\Gamma < 4/3$) no stable core exists and the collapse runs away
regardless of the dissipation (Section~\ref{sec:selfgrav_2d_gpu} and the
role-of-$\gamma$ analysis). The basic, watchable case is therefore the warm,
super-critical cloud, in which collapse, bounce, ring-down and a stable
pressure-supported core all appear cleanly.

\section{Browser-Runnable 3D RealNSE-Gravitation Simulator}\label{sec:cosmology_3d_gpu}

A WebGPU companion simulator
\texttt{cosmology\_3d\_gpu\_100.html} is shipped with the
RealMolecule repository
(\texttt{github.com/Claes542/RealMolecule}). It evolves the
self-gravitating RealNSE system (\ref{eulergrav}) on a $100\times100\times
100$ grid ($10^{6}$ cells) entirely in the browser, using WebGPU
compute shaders for the fluid update and a sub-stepped red-black
Gauss-Seidel relaxation for the Poisson equation $\Delta\varphi = G(\rho - \bar\rho)$.

\paragraph{Initial condition.}
A single Gaussian overdensity at the box centre, both denser and
hotter than the surrounding bath. This is the Big Bang initial
condition of Chapter~\ref{ch:bigbang}: a concentrated hot dense seed
in an otherwise uniform low-density background, with no initial
velocity. The Jeans-swindle source $G(\rho - \bar\rho)$ keeps the
background force-free so that only the density contrast pulls.

\paragraph{Live sliders.}
The simulator exposes the following live controls:
\begin{itemize}
\item \emph{$\gamma$} --- EOS parameter $p = \gamma e = \gamma\rho T$.
\item \emph{Gravity $G$} --- coupling strength.
\item \emph{Seed amp} --- peak overdensity (applied on Reset).
\item \emph{Shock visc $C$} --- shock-viscosity coefficient.
\item \emph{Art visc $\nu$} --- background artificial viscosity.
\item \emph{Seed every} --- particle-trace reset interval.
\item \emph{Trace speed} --- velocity magnification for particle advection.
\item \emph{Fluid substeps/frame} --- numerical resolution in time.
\item \emph{Poisson iters/fluid step} --- gravity convergence.
\end{itemize}

\paragraph{Recommended Big Bang $\to$ Big Crunch settings.}
For a single clean BB$\to$BC cycle followed by a stable virial state:
$G = 500$, seed amp $= 0.5$, $\nu = 1.0$, shock $C = 3.0$, $\gamma = 1$,
Poisson iters $= 10$ or higher for long-time stability.

\paragraph{Reading the display.}
\begin{itemize}
\item \emph{Red particle traces} on the mid-plane: persistent advection
of test particles, redrawn periodically (see \emph{seed every} slider).
During BB the traces radiate outward; during BC they sweep inward;
once virial equilibrium is reached they continue to circulate along
divergence-free streamlines, reflecting the fact that
$\partial_t\rho = 0$ does \emph{not} mean $u = 0$.
\item \emph{Blue polyline} --- density profile $\rho(x_1, x_2^{\mathrm{mid}}, x_3^{\mathrm{mid}})-1$ on the mid-line cross-cut, auto-scaled.
\item \emph{Green polyline} --- temperature profile $T-1$ on the same line.
\item \emph{Cyan polyline} --- gravitational potential $\varphi$ on the same line.
\item \emph{Diagnostic readout} --- peak values of $|\varphi|$, $|\rho-1|$, $|u|$, and simulation time.
\end{itemize}

The simulator is the working companion to the present chapter and
Chapter~\ref{ch:bigbang}: the reader can vary $G$ and the seed amplitude
to scan the Jeans-criterion threshold, watch the
expansion-cooling and compression-warming phases of one Big Bang
followed by one Big Crunch directly in the temperature line plot, and
verify that the cumulative dissipation $D = \int\!\int\mu|\nabla u|^2\,dx\,dt$
prevents the system from oscillating indefinitely.

\section{Browser-Runnable 2D Self-Gravitating Collapse}\label{sec:selfgrav_2d_gpu}

A second WebGPU companion, \texttt{selfgrav\_collapse\_2d\_gpu.html}, reduces
the setting to its essence: a single \emph{isolated} self-gravitating gas cloud
in two dimensions, with the cosmological background and periodicity removed. The
potential satisfies $\Delta\varphi = G\rho$ with the Dirichlet condition
$\varphi = 0$ on the box --- an isolated cloud in empty space rather than a
periodic universe --- solved by Jacobi relaxation; the fluid obeys the same
compressible RealNSE system, the momentum equation carrying the gravitational body
force $-\rho\nabla\varphi$, a stabilizing artificial viscosity $\nu = h$, and the
irreversible dissipation $\nu|\nabla u|^2$ added explicitly to the
internal-energy balance.

\paragraph{Pressureless initial cloud.}
The default initial state is a cold dense disc with internal energy
$e \approx 0$, hence pressure $p = \gamma e \approx 0$: a \emph{pressureless}
(``dust'') cloud at rest. There is then no initial Jeans barrier --- the cloud
free-falls under its own gravity, and \emph{all} of the pressure and temperature
that eventually halt the collapse are generated \emph{during} the collapse, by
compression ($-p\,\nabla\!\cdot u$) and by the dissipation $\nu|\nabla u|^2$.

\paragraph{Two critical relations, and a dimensional caveat.}
The simulator makes the \emph{Jeans threshold} for the onset of collapse visible:
a stiffer (higher-$\gamma$) gas resists with more pressure and needs more gravity,
so lowering the $G$ slider until the cloud no longer contracts locates the
critical coupling. The \emph{outcome} --- whether the collapse halts or runs away
--- is governed by the adiabatic index against the dimension-dependent threshold
$\Gamma_{\rm crit} = 2 - 2/D$. Here lies an essential difference between two and
three dimensions: in 3D, $\Gamma_{\rm crit} = 4/3$ (collapse halts only for
$\gamma > 1/3$), whereas in \emph{this 2D} simulation $\Gamma_{\rm crit} = 1$, so
\emph{any} $\gamma > 0$ halts at a pressure-supported core. The difference is
rooted in the gravity law: $\Delta\varphi = G\rho$ gives the Newtonian
$1/r^2$ force in 3D but a logarithmic, $1/r$ force in 2D --- the 2D run is
physically an infinite gas \emph{filament}, not a sphere. The 2D simulator is
therefore faithful for the collapse \emph{dynamics} and the Jeans onset; the
$4/3$ runaway threshold itself must be studied in the 3D simulator of
Section~\ref{sec:cosmology_3d_gpu}.

\paragraph{Collapse, bounce and ring-down.}
With the pressureless start the cloud first contracts in near free-fall; the
centre runs away to high density and a dense core forms. As the core heats, the
rising pressure halts the infall, but the inward momentum overshoots hydrostatic
balance, so the core \emph{bounces} and then \emph{rings down} through damped
oscillations to a virialized equilibrium. The damping is supplied by
$\nu|\nabla u|^2$: each bounce carries steep velocity gradients whose kinetic
energy is converted irreversibly into heat, so the oscillations decay and the
core settles. The dense core thus forms not by a single smooth contraction but
by overshoot followed by a damped ring-down.

\paragraph{Four cross-cut line plots.}
Along the centre row the simulator draws four profiles on a \emph{logarithmic}
scale (about ten octaves, so the curves keep climbing through a collapse rather
than saturating): density (cyan), temperature (orange), pressure (yellow), and
the gravitational well depth $|\varphi|$ (magenta). The density and pressure
curves spike at the centre, overshoot, dip, and ring down; the potential well
deepens; the temperature climbs from the compression and dissipation heating.

\paragraph{Three-dimensional companion.}
A fully three-dimensional version, \texttt{selfgrav\_collapse\_3d\_gpu.html},
runs the identical model on a $100^3$ grid with the seven-point Poisson solve and
the $1/r^2$ Newtonian force, rendering a movable $z$-slice with the same four
log-scale cross-cuts. It is the simulator in which the genuine
$\Gamma_{\rm crit}=4/3$ threshold lives: with the pressureless dust sphere it
free-falls, forms a dense core, bounces, and rings down for $\gamma>1/3$, while
$\gamma<1/3$ gives runaway --- the behaviour the 2D filament cannot show.

\paragraph{What halts the collapse --- and what dissipation does not do.}
It is tempting to suppose that, with no radiative cooling, the irreversible
heating $\nu|\nabla u|^2$ traps enough energy to halt the collapse for any
$\gamma$. This is \emph{not} so in general, and the distinction is instructive.
What sets whether the collapse halts is the adiabatic index against the
threshold $\Gamma_{\rm crit} = 2-2/D$ of the preceding paragraph: a
\emph{super-critical} gas ($\Gamma > \Gamma_{\rm crit}$) builds pressure faster
than gravity demands and settles into a pressure-supported core, while a
\emph{sub-critical} gas ($\Gamma < \Gamma_{\rm crit}$) does not. The dissipation
$\nu|\nabla u|^2$ supplies the \emph{damping} that lets the bounce ring down to
equilibrium and the entropy production of the 2nd Law --- but it does not by
itself lift a sub-critical gas to stability. Energy conservation explains why:
with no cooling the dissipated infall energy gives an internal energy
$U \sim |W| \sim GM^2/R \propto R^{-1}$, which is \emph{exactly} the marginal
$\Gamma_{\rm crit}$ scaling --- dissipation drives the system \emph{onto} the
critical knife-edge, not above it.

The consequence is dimensional. In \emph{this 2D} run $\Gamma_{\rm crit}=1$, so
\emph{every} $\gamma>0$ is super-critical and the cloud halts at a core --- which
is why the 2D simulator always forms a stable core. In \emph{3D},
$\Gamma_{\rm crit}=4/3$: for $\gamma<1/3$ the gas is sub-critical and
\emph{collapses anyway}, dissipation heating notwithstanding (the 3D simulator's
readout flips to runaway there). Physically this is correct: the protostellar
first core halts because the opaque, adiabatic gas has $\Gamma>4/3$, not because
of turbulent dissipation; and real collapse resumes --- toward a star, and
ultimately, with General Relativity, a \emph{black hole} --- precisely when
endothermic microphysics (molecular dissociation, ionization) drags the effective
$\Gamma$ back below $4/3$, the same heat now absorbed by the phase change rather
than building pressure. What this Newtonian model cannot produce is the black
hole itself: a horizon requires General Relativity (a Schwarzschild radius, the
Chandrasekhar/TOV mass limits), absent here. The model thus marks cleanly where
this level of the hierarchy ends and General Relativity must take over.

\section{Browser-Runnable 3D $\pm$Mass Cosmology: the Cosmic Web as Phase Separation}\label{sec:pmmass_web}

The simulators above keep the density positive, $\rho\ge0$, and follow how a single positive gas gathers into
cores. A further companion, \texttt{cosmology\_darkenergy\_3d\_gpu.html}, removes exactly that constraint and asks
what happens when the mass density is allowed to take \emph{both} signs --- the ``twin Universe with negative
masses'' of \S\,\textit{Gravitons?} made concrete, and the setting in which a \emph{cosmic web} with genuine
\emph{voids} appears.

\paragraph{The $\pm$mass equations.} We keep the self-gravitating RealNSE system (\ref{eulergrav}) unchanged in
\emph{form} but drop the sign constraint on the source: the density $\rho=\Delta\varphi$ is now \emph{two-valued},
taking either sign, while the internal energy is kept positive by the closure $e=|\rho|T$. Writing $u=m/\rho$ for the
velocity and $p=\gamma_s e\ge0$ for the pressure, the system solved is
\begin{equation}\label{pmflowbook}
\begin{array}{rcl}
\Delta\dot\varphi+\nabla\cdot m &=& 0,\\[3pt]
\dot m+\nabla\cdot (mu)+\nabla p-\rho\nabla\varphi &=& \nu\,\Delta u,\\[3pt]
\dot e+\nabla\cdot (eu)+p\,\nabla\cdot u &=& \nu_{\rm sh}\,(\nabla\cdot u)^2,
\end{array}
\end{equation}
with $\rho=\Delta\varphi\gtrless0$ and $e=|\rho|T$. The \emph{only} change from (\ref{eulergrav}) is the two sign
possibilities for $\rho$: the gravitational body force $-\rho\nabla\varphi$ now carries the sign of the local mass,
so \textbf{like mass attracts like and unlike mass repels} --- positive attracts positive, negative attracts
negative, and the two sectors push apart. Because $\rho$ crosses zero the bare velocity $u=m/\rho$ would be singular
there; we hold $|\rho|\ge\rho_{\min}$ (a thin gap the field is pushed out of), so the $+$ and $-$ phases meet at a
\emph{contact wall} rather than passing through $\rho=0$. The shock viscosity $\nu_{\rm sh}\sim h\,|\nabla\cdot u|$
supplies, as before, the entropy production by which pressure arrests collapse, and $|\rho|,|m|,e$ are bounded so
that like-attraction \emph{saturates} at a finite web instead of running to a point singularity.

\paragraph{The cosmic web is a phase separation.} System (\ref{pmflowbook}) is, mathematically, a \emph{bistable
phase separation}. The field $\rho$ has two stable phases ($\rho\gtrless0$) separated by an unstable state at
$\rho=0$ --- a double well, exactly the structure of the Allen--Cahn and Cahn--Hilliard equations of pattern
formation. The segregating force here is not local diffusion but the long-range gravitational kernel
$\varphi=\Delta^{-1}\rho$, so (\ref{pmflowbook}) is a \emph{nonlocal} Cahn--Hilliard of Ohta--Kawasaki type (a
Coulomb kernel coupled to a double well), but with the kernel \emph{attractive} for like mass rather than repulsive:
where the Ohta--Kawasaki repulsion \emph{frustrates} coarsening into finite-scale patterns, gravity's
like-attraction \emph{drives} it, toward hierarchical clustering. Mass is conserved by the flow
($\dot\rho+\nabla\cdot m=0$), so the dynamics is a conserved, momentum-carrying (``model-H'') Cahn--Hilliard. The
upshot: the \textbf{cosmic web is a spinodal-decomposition pattern} --- its filaments, nodes and voids are the
coarsening domain morphology of a bistable phase separation, arrested at finite scale by pressure and dissipation.
This is the \emph{increasing granularity} of \S\,\textit{Cosmology}, now named as the phase-separation process it is.

\paragraph{Reading the display --- and what a ``void'' is.} The simulator seeds a smoothed zero-mean $\pm\rho$ field
(a live \emph{web-scale} slider sets its correlation length), settles $\varphi$ by red-black Gauss--Seidel, and
evolves (\ref{pmflowbook}), drawing either the mid-slice or a line-of-sight projection: red $=$ positive-mass web
(filaments and nodes), blue $=$ negative-mass voids. An \emph{Analyze structure} button classifies the overdense
$+$mass morphology by the eigenvalues of the density Hessian (nodes / filaments / walls) and reports the void
fraction.

A word on that void fraction, since the simulations sharpen a real issue. \emph{Voids are the negative-mass
regions}, so the physical void fraction is the volume fraction with $\rho<0$, and for a zero-mean $\pm$seed it is
$\approx50\%$ at $t=0$ \emph{by construction}. One would like it to \emph{grow} as the positive mass clumps into
thin dense filaments and the negative mass spreads into volume-filling voids --- the ``voids expand'' picture of
dark energy, real cosmic voids occupying most of the volume. In its original signed-inertia form the compressible
run of (\ref{pmflowbook}) does \emph{not} deliver this. Instead the negative phase is progressively consumed and the measured void fraction
\emph{falls} to a few percent (${\sim}2\%$, found consistently in the WebGPU solver and in an independent
finite-difference port), the positive phase saturating to fill the box. The one-sign outcome is independent of the random seed, so it is
structural, not spontaneous symmetry breaking --- and its origin is instructive. The velocity is formed as
$u=m/\rho$ with \emph{signed} $\rho$, which endows the negative-mass regions with negative \emph{inertia}: their
momentum opposes their velocity, $\dot u=F/\rho$ with $\rho<0$, the classic negative-mass \emph{runaway}. The two
phases are therefore \emph{not} mirror images --- the positive phase is an ordinary pressure-supported gas, the
negative phase a runaway-unstable fluid --- so the negative phase cannot self-support, is consumed, and the voids
collapse (the hard $|\rho|$ caps, which break mass conservation once hit, merely amplify the winner). This also
departs from the \emph{formal} $\pm$mass of the analysis above, where the sign lives in the curvature of $\varphi$
and the energy/inertia stays positive. Recovering a symmetric, void-dominated web thus requires giving \emph{both}
phases positive inertia --- taking the velocity from $|\rho|$, $u=m/(|\rho|+\rho_{\min})$, with the sign of $\rho$
entering \emph{only} the gravitational coupling --- together with a mass-conserving transport in place of hard
clipping. Then $+$ and $-$ mass are identical ordinary gases that each self-attract and mutually repel, and
segregate into a $\sim\!50/50$ $\pm$web; a reduced, near-conservative semi-Lagrangian model (positive-inertia
backtrace advection by $u=-\operatorname{sign}(\rho)\nabla\psi$, $\Delta\psi=\rho$) already realises this and its
$\sim\!50\%$ voids, and is the cleaner demonstration of the morphology. The WebGPU solver has since been switched
to this \emph{Route~A} scheme --- positive inertia $u=m/(|\rho|+\rho_{\min})$, continuity transporting the signed
charge with the conservative flux $\rho u$, and no signed amplitude cap --- and an offline finite-difference port
confirms the void fraction then \emph{holds} near $50\%$ with total charge conserved and $|\rho|$ bounded, the
symmetric $\pm$web restored. What remains open is \emph{sharpness}: at settings safely inside the stability region
the centered advection is over-diffusive, so the web stays smooth (only mild clumping); a crisp, filamentary web at
fixed stability calls for an upwind or semi-Lagrangian transport, the natural next step. (One further, benign diagnostic --- the fraction of the \emph{smoothed positive} field below
$0.3\times$ its mean --- also reads a few percent, but for an unrelated reason: after box-smoothing only the deep
interiors of the largest holes survive that threshold. It measures deep $+$underdensity, not the negative-mass
volume, and the two must not be confused.)

\paragraph{$\pm$mass as \emph{contrast}, not substance.} A remark on how this section stands to the article
\emph{Coulomb $+$ Newton: Towards a Cosmology}, which reads the same two-valued Poisson equation differently. There
the gravitational source is not a signed substance but the \emph{departure of an everywhere-positive energy density
from its own mean},
\begin{equation}\label{eq:contrastsource}
\Delta\varphi_g \;=\; 4\pi G\,(\rho-\bar\rho),
\end{equation}
so the source takes both signs while the matter does not: ``negative mass'' names an \emph{underdense region},
nothing more. Mathematically this is the same solver --- $\rho-\bar\rho$ is precisely the zero-mean signed field
seeded above --- and the simulations therefore carry over unchanged, with red read as overdense and blue as
underdense. What the contrast reading buys is exactly the difficulty diagnosed in the previous paragraph: with
$\rho\ge0$ everywhere, inertia is positive everywhere by construction, the negative-mass runaway cannot arise, and
\emph{Route~A} --- taking the velocity from $|\rho|$ and letting the sign enter only the gravitational coupling ---
stops being a repair and becomes the natural discretisation. The price is that voids no longer repel by carrying
negative mass but simply fail to attract, so the outward push on their surroundings is relative, not absolute; the
article accordingly claims a coasting expansion ($q=0$) rather than an accelerating one, and does not offer the
$\pm$web as dark energy. Both readings are kept here: the signed-substance form because it is what the solver
literally integrates and because its failure mode is instructive, the contrast form because it is what the article
defends.



